\subsection*{选择重放哪些内容}

恢复是否一致，需要先规定比较哪些可观察字段。一份响应中的 P1、终态、O1、D1、Q7、币种、3000 与当次 7000 构成这里的业务结果投影，同一 S 的各次同参恢复应当保持这组字段。Date 是响应生成时间字段，Request-Id 是本次请求编号；它们属于通信投影，由本次响应填写，分别表示何时回答、回答哪次尝试，允许变化。

以首次成功与十分钟后恢复为例，两份响应的业务投影都含 P1、O1、D1、3000 和 7000；通信投影分别含 R1 与 R2。若恢复把 7000 换成当前 6000，即使其他字段不变，业务投影也已不同。若仅把 R1 换成 R2，则仍满足原结果规则。这项观察范围直接决定测试该比什么，也使接口可以在保留诊断能力的同时对业务内容作出可核对的稳定承诺。

\begin{figure}[htbp]
\centering
\begin{tikzpicture}[x=1cm,y=1cm]
\node[bookresult,text width=10cm] (b) at (0,0) {共同业务投影：P1 / O1 / D1 / CNY 3000 / 当次余额 7000};
\node[booknode,text width=4cm] (r1) at (-3,-1.8) {首次响应\\201 / R1 / 首次 Date};
\node[booknode,text width=4cm] (r2) at (3,-1.8) {恢复响应\\200 / R2 / 本次 Date};
\draw[bookarrow] (b)--(r1);
\draw[bookarrow] (b)--(r2);
\node[booknote,text width=10cm] at (0,-3.1) {通信投影随尝试变化\\把业务余额 7000 替换成当前 6000 则改变原结论};
\end{tikzpicture}
\caption{比较同参终态恢复的业务投影，无须要求整个响应逐字节相同。}
\label{fig:ch05-result-projection}
\end{figure}


可以据此精确说明同一操作的两个响应怎样相等。用 \(r_1,r_2\) 表示两次已经读到终态的同参响应，用 \(\pi_B\) 表示选出上述业务字段的操作，\(\pi_C\) 表示选出本次通信字段的操作。所要求的是
\[
 \pi_B(r_1)=\pi_B(r_2)=R(S).
\]
这里并未要求 \(\pi_C(r_1)=\pi_C(r_2)\)。以数值核对，\(\pi_B\) 中的扣款额都是 3000，当次余额都是 7000；\(\pi_C\) 中的请求编号分别是 R1、R2。这个等式只比较成功恢复出来的结果，不把尚无终态的 404、暂时不可用的 503 或异参的 409 也要求成同一成功响应。

本章把恢复旧业务内容称为结果重放：服务从持久完成记录取出原结果，再为本次请求构造响应。这里选择业务内容稳定，HTTP 包装按本次通信重新产生。原购买编号、终态、订单、扣款、成交快照、当次余额和结果版本保持一致；Date、Request-Id 与恢复标志对应本次尝试。此处响应实际上不一定逐字节一致。

逐字节回放是另一种要求：首次响应的状态行、选定响应头和正文编码要保存为可以原样发送的字节。若签名验证依赖正文的精确编码，或调用者确实需要同一份签名文件，就可能需要这类对象及其长期验证方式。讨论到现在，我们实现的购买接口没有这一要求。JSON 成员排序或无意义空白变化不改变这里保存的业务结果；任何金额、编号、终态或字段语义变化都不能用编码差异解释。

这个选择也影响 HTTP 状态。首次事务明确完成新购买时，本接口返回 201 Created，表示本次请求形成了资源。后来同参 POST 找到已提交终态时返回 200 OK，表示成功提供原结果；GET 查询成功也返回 200。201 与 200 因本次入口动作不同而不同，稳定业务内容仍相同。HTTP 状态码的基本含义来自 RFC 9110，本章的重试规则是购买接口额外规定的行为，HTTP 本身不会自动识别 K1\cite{http}\footnote{本书使用的 Request-Id、Operation-Permit 与 Idempotency-Replayed 为示例接口字段，不能据此假定任意 HTTP 服务已经实现同样的购买恢复协议。}。

\subsection*{一份原结果与两次响应}

首次成功的响应可以写成下面的形式。Location 指向 P1 的购买资源，Content-Type 说明正文使用 JSON；\code{Cache-Control: no-store} 要求遵循规则的缓存不要保存这份含账户信息的响应\cite{httpcache}。Date 是本次响应产生的时间信息。Request-Id 和 Idempotency-Replayed 是本接口采用的诊断头，前者对应当前尝试，后者说明此次是否从已保存终态恢复；它们不参与 S 的唯一性判断。

\begin{BookCode}
HTTP/1.1 201 Created
Location: /v1/purchases/P1
Content-Type: application/json
Cache-Control: no-store
Date: Sat, 17 Oct 2026 10:00:01 GMT
Request-Id: R1
Idempotency-Replayed: false

{"result_version":"purchase-result-v1",
 "purchase_id":"P1","status":"succeeded",
 "order_id":"O1","debit_id":"D1",
 "product_id":"paper-001","quantity":1,
 "quote_id":"Q7","currency":"CNY",
 "amount_minor":3000,"balance_after_minor":7000}
\end{BookCode}

假设这份响应丢失，R2 十分钟后找到 P1。其响应为 200 OK，Request-Id 改为 R2，Date 表示十分钟后的响应时点，Idempotency-Replayed 为 true，正文仍取上面保存的 \code{result\_payload}。R2 不把自己到达的时间写入 \code{result\_recorded\_at}，也不把当前账户 6000 填入 \code{balance\_after\_minor}。服务重放的是 P1 的业务内容，记录的是 R2 这一次响应经过。

\begin{figure}[htbp]
\centering
\begin{tikzpicture}[x=1cm,y=1cm]
\node[bookresult,text width=7cm] (p) at (0,0) {保存的原结果：P1 / O1 / D1 / 成功\\CNY 3000 / 当次余额 7000};
\node[booknode,text width=4.5cm] (a) at (-3,-2) {首次响应：201\\R1 / 10:00:01\\Replayed=false};
\node[booknode,text width=4.5cm] (b) at (3,-2) {恢复响应：200\\R2 / 10:10:01\\Replayed=true};
\draw[bookarrow] (p)--(a);
\draw[bookarrow] (p)--(b);
\end{tikzpicture}
\caption{两个响应包装同一份业务结果，各自记录本次通信尝试。}
\label{fig:ch05-response-attempts}
\end{figure}


若首次事务提交后、响应发送前进程退出，客户端从未看到 R1 的响应字节，R2 仍能得到同样的业务结果。这说明本例需要保存的是完成事实，不能只依靠网关曾经缓存过什么响应。把结果保存到内存响应缓存，数据库提交与缓存写入之间仍有故障窗口；缓存失效或进程重启也会使原内容消失。缓存可加速读取，缺失时应回到持久完成记录。

同样，不必把所有 HTTP 头都保存为业务历史。连接管理、压缩方式、访问凭据续期之类信息服务于当前连接或当前安全状态，旧头原样发送可能已经不适用。重放规则应明确列出稳定内容和每次产生的包装。若将来某个头包含成交凭证的签名或结果版本，它与历史结果有实质关系，就需要纳入保存和兼容设计，不能一律当作可随意变化的头。

\subsection*{版本升级以后仍能解释原结果}

\code{result\_version} 表示这份原结果按哪套字段规则解释。版本 v1 规定 \code{amount\_minor} 与 \code{balance\_after\_minor} 都是整数分，\code{status=succeeded} 表示创建购买与本地扣款成功。服务升级到 v2，加入预计送达时间时，旧结果没有这一承诺。恢复 v1 仍可按 v1 输出，或者显式提供兼容表示，说明新增字段只是本次查询资料。直接用当前代码把旧 payload 解为 v2，再填一个看似合理的预计日期，会给过去的购买增加从未形成的事实。

本章选择在恢复期限内保留 v1 的解码与输出能力。服务找到旧记录后按 \code{result\_version} 选择响应构造器，构造器只使用保存字段，不重新执行计价、库存或余额判断。新操作默认采用新规则，不影响已经绑定的 S。动作版本、输入版本与结果版本可以分别推进，部署前需要检查旧组合仍有完整恢复路径。

若旧结果版本已无法输出，服务应给出明确的兼容性错误并保留业务事实查询，不重新执行购买。数据库有 P1 并不意味着任意版本代码都能正确解释它；记录与解释能力共同决定结果可恢复。恢复期限内撤掉旧解释器，会使服务主动破坏自己承诺的恢复范围。超过期限后如何保留较小的业务证据，则与记录回收策略一起决定。

\begin{figure}[htbp]
\centering
\begin{tikzpicture}[x=1cm,y=1cm]
\node[booknode,text width=3cm] (a) at (0,0) {旧记录 result-v1\\P1 / D1 / 7000};
\node[booknode,text width=3cm] (b) at (4,0) {当前服务\\按版本读取与映射};
\node[bookresult,text width=3cm] (c) at (8,0) {兼容响应\\仍表达 P1 / D1 / 7000};
\draw[bookarrow] (a)--(b);
\draw[bookarrow] (b)--(c);
\node[booknote,text width=10cm] at (4,-1.6) {字段可以有明确的兼容映射；当前余额 6000 不能替换原扣款后的 7000};
\end{tikzpicture}
\caption{结果版本帮助新服务继续解释旧记录，兼容映射保留原业务含义。}
\label{fig:ch05-result-versions}
\end{figure}


版本兼容也要回到业务投影检查。把旧 v1 映射到兼容表示时，映射应保留 P1、D1、CNY 3000 与当次 7000 的含义；金额单位转换可以改变显示字符串，却应明确转换规则和版本。若兼容表示补充当前账户 6000，应把它放入本次查询字段，不能替换原当次余额。由此检查升级，不只检查新 JSON 能否解析，还检查旧购买在升级前后的业务结论是否相同。

字段脱敏也要区分观察权限和原结果。若结果含收货地址等敏感内容，后来的调用者权限可能只允许看到部分字段，服务可提供受限视图并明确其含义，不能把原存储内容直接清空后声称仍能精确恢复。本例结果只列购买与资金标识，但仍按账户权限访问。稳定结果并不取消每次访问的认证与授权检查。

\subsection*{把业务关联接进诊断链路}

请求 ID 标识一次通信尝试，R1、R2、R3 应能分别检索。分布式追踪把一次请求经过的网关、服务和数据库调用组织成相互关联的处理片段，帮助定位耗时或错误；其 \code{trace\_id} 标识这一组诊断片段。不同重试可以属于不同追踪，诊断资料也可能采样或到期删除，因此业务恢复仍从完整 S 出发。

一条诊断事件可同时记录 R2、当前 \code{trace\_id}、S、查到的 P1、恢复来源和返回业务码。候选 P2 若未实际插入，应标明 \code{candidate\_only}；否则运维人员看到 P1、P2 两个号，容易把编号分配次数数成购买次数。记录一次重放事件也不修改 P1 的原终态。业务事实保留多久与日志保留多久，可以有各自策略，恢复不得依赖所有旧请求日志都还存在。

服务端先提交 P1，随后记录发送成功，但连接在客户端接收以前断开，日志也只能证明服务做过发送步骤。客户端真正收到完整响应，还可能在本地保存结果前退出。若产品需要追踪客户端是否确认，应增加明确确认动作和自己的身份规则。它不能从一个响应发送日志自动推出，也不能仅凭确认缺失就重新扣款。本章始终以持久业务记录解释购买，以请求日志解释尝试经过。
\FloatBarrier
